\section{Bell Inequalities and \texorpdfstring{$\tau_n$}{tau_n}}
\paragraph*{Theorem 8.1 (Bell consistency).}
The proper time $\tau_n$ is not a local hidden variable in Bell's sense. The framework is consistent with quantum-mechanical violation of Bell inequalities.

\paragraph*{Proof.}
Bell's theorem assumes a hidden variable $\lambda$ assigned independently to each particle, with measurement outcomes $A(\lambda_A,a)$ and $B(\lambda_B,b)$ and statistically independent variables $\lambda_A,\lambda_B$. In the present framework, however, $\tau_n$ belongs to the pair $(f,Jf)$ rather than to each element separately; one has
\[
\lambda_A = \lambda_B = \tau_n.
\]
Thus $\tau_n$ is a contextual variable, and Bell's theorem does not exclude violation by such theories. The correlation function $E(a,b)$ follows from the structure of $H_+^{\,+}$: since $C(f) \in H_+^{\,+}$ carries an irreducible representation of the Weyl algebra $W(A_1)$ (Stone--von Neumann, [1]), the standard quantum correlation
\[
E(a,b)=\cos \theta_{ab}
\]
is reproduced, yielding the CHSH value
\[
S = 2\sqrt{2} > 2.
\]
Hence $\tau_n$ explains why the correlation exists (shared proper time fixed at the near-collision), while Stone--von Neumann determines the strength of the violation.

\paragraph*{Remark (What $\tau_n$ does and does not do).}
The quantity $\tau_n$ encodes the magnitude $|z_-|$ of the displacement from the photon sector at the moment of the near-collision. It does not determine the measurement direction. When $f$ is measured, $z_-$ is projected onto a value; the paired state $Jf$ is projected onto $-z_-$, the $J$-image. The anticorrelation is built in at the collision event and is not transmitted at measurement. No superluminal signal is required.

